\documentclass[10pt, letterpaper]{article}

% Packages:
\usepackage[
    ignoreheadfoot, % set margins without considering header and footer
    top=2 cm, % seperation between body and page edge from the top
    bottom=2 cm, % seperation between body and page edge from the bottom
    left=2 cm, % seperation between body and page edge from the left
    right=2 cm, % seperation between body and page edge from the right
    footskip=1.0 cm, % seperation between body and footer
    % showframe % for debugging 
]{geometry} % for adjusting page geometry
\usepackage{titlesec} % for customizing section titles
\usepackage{tabularx} % for making tables with fixed width columns
\usepackage{array} % tabularx requires this
\usepackage[dvipsnames]{xcolor} % for coloring text
\definecolor{primaryColor}{RGB}{0, 79, 144} % define primary color
\usepackage{enumitem} % for customizing lists
\usepackage{fontawesome5} % for using icons
\usepackage{amsmath} % for math
\usepackage[
    pdftitle={Tosin's CV},
    pdfauthor={Leonardo Tosin},
    pdfcreator={LaTeX with RenderCV},
    colorlinks=true,
    urlcolor=primaryColor
]{hyperref} % for links, metadata and bookmarks
\usepackage[pscoord]{eso-pic} % for floating text on the page
\usepackage{calc} % for calculating lengths
\usepackage{bookmark} % for bookmarks
\usepackage{lastpage} % for getting the total number of pages
\usepackage{changepage} % for one column entries (adjustwidth environment)
\usepackage{paracol} % for two and three column entries
\usepackage{ifthen} % for conditional statements
\usepackage{needspace} % for avoiding page brake right after the section title
\usepackage{iftex} % check if engine is pdflatex, xetex or luatex

% Ensure that generate pdf is machine readable/ATS parsable:
\ifPDFTeX
    \input{glyphtounicode}
    \pdfgentounicode=1
    % \usepackage[T1]{fontenc} % this breaks sb2nov
    \usepackage[utf8]{inputenc}
    \usepackage{lmodern}
\fi



% Some settings:
\AtBeginEnvironment{adjustwidth}{\partopsep0pt} % remove space before adjustwidth environment
\pagestyle{empty} % no header or footer
\setcounter{secnumdepth}{0} % no section numbering
\setlength{\parindent}{0pt} % no indentation
\setlength{\topskip}{0pt} % no top skip
\setlength{\columnsep}{0cm} % set column seperation
\makeatletter
\let\ps@customFooterStyle\ps@plain % Copy the plain style to customFooterStyle
\patchcmd{\ps@customFooterStyle}{\thepage}{
    \color{gray}\textit{\small Leonardo Tosin - Page \thepage{} of \pageref*{LastPage}}
}{}{} % replace number by desired string
\makeatother
\pagestyle{customFooterStyle}

\titleformat{\section}{\needspace{4\baselineskip}\bfseries\large}{}{0pt}{}[\vspace{1pt}\titlerule]

\titlespacing{\section}{
    % left space:
    -1pt
}{
    % top space:
    0.3 cm
}{
    % bottom space:
    0.2 cm
} % section title spacing

\renewcommand\labelitemi{$\circ$} % custom bullet points
\newenvironment{highlights}{
    \begin{itemize}[
        topsep=0.10 cm,
        parsep=0.10 cm,
        partopsep=0pt,
        itemsep=0pt,
        leftmargin=0.4 cm + 10pt
    ]
}{
    \end{itemize}
} % new environment for highlights

\newenvironment{highlightsforbulletentries}{
    \begin{itemize}[
        topsep=0.10 cm,
        parsep=0.10 cm,
        partopsep=0pt,
        itemsep=0pt,
        leftmargin=10pt
    ]
}{
    \end{itemize}
} % new environment for highlights for bullet entries


\newenvironment{onecolentry}{
    \begin{adjustwidth}{
        0.2 cm + 0.00001 cm
    }{
        0.2 cm + 0.00001 cm
    }
}{
    \end{adjustwidth}
} % new environment for one column entries

\newenvironment{twocolentry}[2][]{
    \onecolentry
    \def\secondColumn{#2}
    \setcolumnwidth{\fill, 4.5 cm}
    \begin{paracol}{2}
}{
    \switchcolumn \raggedleft \secondColumn
    \end{paracol}
    \endonecolentry
} % new environment for two column entries

\newenvironment{header}{
    \setlength{\topsep}{0pt}\par\kern\topsep\centering\linespread{1.5}
}{
    \par\kern\topsep
} % new environment for the header

\newcommand{\placelastupdatedtext}{% \placetextbox{<horizontal pos>}{<vertical pos>}{<stuff>}
  \AddToShipoutPictureFG*{% Add <stuff> to current page foreground
    \put(
        \LenToUnit{\paperwidth-2 cm-0.2 cm+0.05cm},
        \LenToUnit{\paperheight-1.0 cm}
    ){\vtop{{\null}\makebox[0pt][c]{
        \small\color{gray}\textit{Última atualização em Junho 2025}\hspace{\widthof{Last updated in September 2024}}
    }}}%
  }%
}%

% save the original href command in a new command:
\let\hrefWithoutArrow\href

% new command for external links:
\renewcommand{\href}[2]{\hrefWithoutArrow{#1}{\ifthenelse{\equal{#2}{}}{ }{#2 }\raisebox{.15ex}{\footnotesize \faExternalLink*}}}


\begin{document}
    \newcommand{\AND}{\unskip
        \cleaders\copy\ANDbox\hskip\wd\ANDbox
        \ignorespaces
    }
    \newsavebox\ANDbox
    \sbox\ANDbox{}

    \placelastupdatedtext
    \begin{header}
        \textbf{\fontsize{24 pt}{24 pt}\selectfont Leonardo Tosin}

        \vspace{0.3 cm}

        \normalsize
        \mbox{{\color{black}\footnotesize\faMapMarker*}\hspace*{0.13cm}Garibaldi, RS - Brazil}%
        \kern 0.25 cm%
        \AND%
        \kern 0.25 cm%
        \mbox{\hrefWithoutArrow{leob.tosin@hotmail.com}{\color{black}{\footnotesize\faEnvelope[regular]}\hspace*{0.13cm}leob.tosin@hotmail.com}}%
        \kern 0.25 cm%
        \AND%
        \kern 0.25 cm%
        \mbox{\hrefWithoutArrow{tel:+55 (54) 99929-9478}{\color{black}{\footnotesize\faPhone*}\hspace*{0.13cm}+55 (54) 99929-9478}}%
        \kern 0.25 cm%
        \AND%
        \kern 0.25 cm%
        \mbox{\hrefWithoutArrow{https://leotosin.com.br/}{\color{black}{\footnotesize\faLink}\hspace*{0.13cm}leotosin.com.br}}%
        \kern 0.25 cm%
        \AND%
        \kern 0.25 cm%
        \mbox{\hrefWithoutArrow{https://www.linkedin.com/in/leonardo-tosin-b57406112/}{\color{black}{\footnotesize\faLinkedinIn}\hspace*{0.13cm}leonardo-tosin}}%
        \kern 0.25 cm%
        \AND%
        \kern 0.25 cm%
        \mbox{\hrefWithoutArrow{https://github.com/TheLe0}{\color{black}{\footnotesize\faGithub}\hspace*{0.13cm}TheLe0}}%
    \end{header}

    \vspace{0.3 cm - 0.3 cm}


    \section{Sobre}



        
        \begin{onecolentry}
            Iniciei minha carreira em desenvolvimento de software em 2017 e, desde então, atuei continuamente como Engenheiro de Software Full Stack. Ao longo dos anos, adquiri experiência com uma ampla variedade de tecnologias e tipos de produtos, o que me proporcionou uma compreensão sólida de todo o ciclo de vida de desenvolvimento de software.
        
            Minha principal especialidade está em .NET, Java, Angular e plataformas em nuvem como AWS e Azure. Tenho ampla experiência na construção de aplicações de alta performance e baixa latência, capazes de lidar com grandes volumes de dados.
        
            Um dos meus principais interesses está em sistemas de banco de dados — tanto SQL quanto NoSQL. Gosto de me aprofundar nos seus mecanismos internos para entender melhor como funcionam e como otimizá-los para diferentes cenários.
        
            Sou bacharel em Ciência da Computação e, como trabalho de conclusão de curso, desenvolvi uma plataforma de gamificação para roteiros turísticos. Você pode saber mais sobre ela assistindo a este \href{https://www.youtube.com/watch?v=xZLdsME5gGU}{vídeo}.
        \end{onecolentry}

        \vspace{0.2 cm}
    
    \section{Links}
        \begin{onecolentry}
            \begin{highlightsforbulletentries}
                \item \href{https://leotosin.com.br/}{Site pessoal}
                \item \href{https://medium.com/@TheLe0}{Blog}
                \item \href{https://github.com/TheLe0}{GitHub}
                \item \href{https://hub.docker.com/u/thele0}{Docker Hub}
                \item \href{https://www.nuget.org/profiles/TheLe0/}{NuGet}
                \item \href{https://stackoverflow.com/users/9767014/thele0?tab=profile}{Stack Overflow}
            \end{highlightsforbulletentries}
        \end{onecolentry}
    \section{Educação}


        \begin{twocolentry}{
            
            
        \textit{Março 2013 – Agosto 2022}}
            \textbf{Universidade de Caxias do Sul}

            \textit{Bacharelado em Ciência da Computação}
        \end{twocolentry}

        \vspace{0.10 cm}
        \begin{onecolentry}
            \begin{highlights}
                \item Minha nota final foi 9.8/10
                \item Desenvolvi como trabalho de conclusão de curso um aplicativo mobile chamado Itinehapp — ou, em tradução livre, "Seu aplicativo de roteiros felizes" — que permite que turistas montem seus próprios roteiros para destinos ao redor do mundo e compartilhem com outras pessoas.
            \end{highlights}
        \end{onecolentry}



    
    \section{Experiência}

        \begin{twocolentry}{
        \textit{Remoto}    
            
        \textit{Janeiro 2024 - atualmente}}
            \textbf{Senior Full Stack Engineer}
            
            \textit{CI\&T}
        \end{twocolentry}

        \vspace{0.10 cm}
        \begin{onecolentry}
            \begin{highlights}
                \item Atuação com um cliente dos Estados Unidos no desenvolvimento de um sistema ECM (Enterprise Content Management) capaz de fazer upload de milhões de documentos diariamente, com um mecanismo de busca rápido e amigável baseado em classificação por metadados.
                \item Trabalho em todo o ciclo de vida da plataforma, desde implementações e melhorias no front-end, com foco em uma interface mais intuitiva, até otimizações no back-end, garantindo alto desempenho e escalabilidade.
                \item Também contribuo em diversas áreas como Cloud, DevOps, FinOps, SRE, Observabilidade e manutenção/otimização de banco de dados.
                \item Algumas contribuições importantes do meu trabalho:
                \begin{highlightsforbulletentries}
                    \item Participação ativa nos esforços de conformidade para obtenção das certificações SOC 2 e SOC 3, incluindo a implementação de melhorias de segurança solicitadas pelo AWS Security Hub e de uma estratégia de Recuperação de Desastres para failover entre regiões.
                    \item Implementação de SSO (Single Sign-On) e SCIM (System for Cross-domain Identity Management), permitindo integração com provedores de identidade como Active Directory, Okta, Google e Amazon, com sincronização bidirecional de usuários e grupos.
                    \item Desenvolvimento de mecanismo de busca de documentos combinando técnicas de Full Text Search e Busca Semântica, possibilitando pesquisas em subsegundos sobre milhões de documentos.
                    \item Criação de funcionalidades com inteligência artificial, incluindo classificação automática de documentos, análise por base de conhecimento e sumarização de conteúdo.
                    \item Desenvolvimento de conectores robustos para integração com sistemas corporativos como SAP, Salesforce e Box, permitindo integrações multiplataforma.
                \end{highlightsforbulletentries}
                \item Ferramentas utilizadas: C\#, .NET, .NET Framework, Java, MongoDb, AWS, Apache Lucene, Serverless, API Gateway, AWS S3, Vector Search e Python.
            \end{highlights}
        \end{onecolentry}

        \begin{twocolentry}{
        \textit{Remote}    
            
        \textit{Júlio 2022 - Janeiro 2024}}
            \textbf{Engenheiro de Software Pleno}
            
            \textit{Warren Investimentos}
        \end{twocolentry}

        \vspace{0.10 cm}
        \begin{onecolentry}
            \begin{highlights}
                \item Levantamento de requisitos técnicos e funcionais em conjunto com o time para desenvolvimento de novas funcionalidades escaláveis e resilientes para o negócio.
                \item Apoio a novos integrantes da equipe por meio de code reviews, sessões 1:1 e pair programming, facilitando a adaptação ao sistema e ao domínio do negócio.
                \item Desenvolvimento de funcionalidades com foco principal em .NET e NodeJS, sempre com alta cobertura de testes (unitários, integração, ponta a ponta e de estresse).
                \item Integração entre back-end e front-end via BFFs (Backends for Frontends), realizando cruzamento de dados com múltiplos microsserviços.
                \item Criação de aplicações web voltadas ao backoffice, otimizando operações do dia a dia.
                \item Atuação na identificação e resolução de bugs, participação ativa em war rooms para análise e mitigação de incidentes.
                \item Participação em cerimônias ágeis (Scrum), como dailies, plannings, refinamentos teóricos e técnicos, além de retrospectivas de sprint.
                \item Documentação das funcionalidades desenvolvidas e fluxos das aplicações.
                \item Monitoramento e gestão de recursos em nuvem, tanto via portal da plataforma quanto por meio de IaC com Terraform.
                \item Contato direto com stakeholders do produto para entender suas rotinas, dores e necessidades, contribuindo para a entrega de funcionalidades mais intuitivas e eficientes.
                \item Ferramentas utilizadas: C\#, .NET, .NET Framework, MySQL, Apache Kafka, VueJS, NodeJS, MongoDb, AWS, ElasticSearch, DataDog, OpenTelemetry e ECS.
            \end{highlights}
        \end{onecolentry}

        \begin{twocolentry}{
        \textit{Carlos Barbosa, RS}    
            
        \textit{Setembro 2019 – Júlio 2022}}
            \textbf{Desenvolvedor de Software}
            
            \textit{Cooperativa Santa Clara}
        \end{twocolentry}

        \vspace{0.10 cm}
        \begin{onecolentry}
            \begin{highlights}
                \item Correção e aprimoramento do sistema de varejo, responsável por importar pedidos do sistema de PDV para o ERP, realizando a criação e faturamento automático dos pedidos.
                \item Implementação de um processo de liquidação financeira em lote utilizando programação concorrente, permitindo o processamento de centenas de títulos em poucos minutos — tarefa que anteriormente era feita manualmente pelo usuário.
                \item Criação de uma solução de armazenamento de arquivos estáticos na nuvem (imagens, PDFs e textos), utilizando o Azure Blob Storage. A integração foi feita por meio da ferramenta AzCopy da Microsoft, que mapeia um diretório local monitorado por um serviço do Azure que detecta e sincroniza automaticamente alterações como criação, remoção ou modificação de arquivos.
                \item Desenvolvimento de uma API REST serverless utilizando Azure Functions e Azure API Management para autenticação e autorização de usuários via Active Directory.
                \item Desenvolvimento do back-end da plataforma de E-commerce da empresa, com integração via REST, WSDL e SOAP com sistemas internos, gateways de pagamento, transportadoras e outros. Atuei desde o início do projeto, não apenas na implementação, mas também na modelagem de software e definição da arquitetura da solução.
                \item Ferramentas utilizadas: C\#, .NET, .NET Framework, SQL Server, Reporting Services, Azure, Azure Blob Storage, Azure Functions, NodeJS e ReactJS.
            \end{highlights}
        \end{onecolentry}
        
        
        \begin{twocolentry}{
        \textit{Garialdi, RS}    
            
        \textit{Abrul 2017 – Setembro 2019}}
            \textbf{Full Stack Developer}
            
            \textit{NuvolaHost Serviços de TI LTDA}
        \end{twocolentry}

        \vspace{0.10 cm}
        \begin{onecolentry}
            \begin{highlights}
                \item Modernização de layouts de páginas web, com foco na criação de interfaces responsivas (UI responsivas).
                \item Desenvolvimento de funcionalidades e sistemas utilizando TDD (Test-Driven Development) durante o processo de implementação.
                \item Apoio na implementação de práticas DevOps, automatizando e aprimorando todo o fluxo de testes e deploy de novas versões de software. Ênfase na execução de aplicações em Kubernetes e gerenciamento de clusters via OpenShift.
                \item Migração de todos os repositórios da empresa de SVN para Git como sistema de controle de versão.
                \item Criação de uma plataforma SaaS de E-mail Marketing, responsável pelo envio de milhares de e-mails diariamente.
                \item Ferramentas utilizadas: PHP, Zend Framework, MySQL, JQuery, Ruby on Rails, Docker, JavaScript, HTML e CSS.
            \end{highlights}
        \end{onecolentry}


        \vspace{0.2 cm}

        \begin{twocolentry}{
        \textit{Caxias do Sul, RS}    
            
        \textit{Júlio 2014 – Março 2015}}
            \textbf{Iniciação Científica}
            
            \textit{Microsoft}
        \end{twocolentry}

        \vspace{0.10 cm}
        \begin{onecolentry}
            \begin{highlights}
                \item Desenvolvi um projeto de mestrado chamado “Mantas Network”, focado na utilização de mantas poliméricas para extração de vazamentos de óleo, gasolina e derivados em ambientes offshore.
                \item Minha principal função foi realizar pesquisa científica de artigos sobre novos métodos e técnicas aplicáveis, além de executar testes práticos e aplicar algoritmos estatísticos e probabilísticos para avaliar a eficácia das soluções propostas.
                \item Ferramentas utilizadas: Python, R, Matlab, Power BI e Excel.
            \end{highlights}
        \end{onecolentry}



    
    \section{Artigos}

    \begin{onecolentry}
        \begin{highlightsforbulletentries}


        \item \href{https://medium.com/@TheLe0/web-scrapping-with-net-and-selenium-webdriver-d8a888756733}{Web Scrapping with .NET and Selenium Webdriver}
        \item \href{https://medium.com/@TheLe0/cloud-flavours-on-azure-the-different-types-of-cloud-services-6e8a12919d78}{Cloud flavours on Azure: The different types of cloud services}
        \item \href{https://medium.com/@TheLe0/learning-how-to-use-aws-cdk-with-net-as-the-iac-to-your-projects-c9806e5ea739}{Learning how to use AWS CDK with .NET as the IaC to your projects}
        \item \href{https://medium.com/@TheLe0/how-to-easily-build-and-deploy-a-website-on-azure-static-web-apps-32e62861a2e5}{How to easily build and deploy a website on Azure Static Web Apps}

        \end{highlightsforbulletentries}
    \end{onecolentry}

    
    \section{Projetos}


        \begin{twocolentry}{
            
            
        \textit{\href{https://www.berteleadvocacia.com.br/}{berteleadvocacia.com.br}}}
            \textbf{Bertele Advocacia}
        \end{twocolentry}

        \vspace{0.10 cm}
        \begin{onecolentry}
            \begin{highlights}
                \item Website para um escritório de advocacia
                \item Ferramentas usadas: PHP, Blade, Laravel, Azure, Docker, Container Apps e Container Registry.
            \end{highlights}
        \end{onecolentry}

        \begin{twocolentry}{
            
            
        \textit{\href{https://www.edrian.com.br/}{edrian.com.br}}}
            \textbf{Engenheiro Édrian}
        \end{twocolentry}

        \vspace{0.10 cm}
        \begin{onecolentry}
            \begin{highlights}
                \item Website para uma consultoria de Engenharia de Segurança do Trabalho
                \item Ferramentas usadas: Angular 17, Azure Static Web Apps e Mailgun.
            \end{highlights}
        \end{onecolentry}

    
    \section{Tecnologias}

        \begin{onecolentry}
            \textbf{Linguagens:} C\#, Java/Kotlin, JS/TS e PHP
        \end{onecolentry}

        \vspace{0.2 cm}

        \begin{onecolentry}
            \textbf{Frameworks:} .NET, Quarkus e Laravel
        \end{onecolentry}

        \vspace{0.2 cm}

        \begin{onecolentry}
            \textbf{Cloud:} Azure e AWS
        \end{onecolentry}

        \vspace{0.2 cm}
    
        \begin{onecolentry}
            \textbf{Banco de Dados:} SQL Server, MongoDb, CosmosDb, Redis, MySQL e PostgreSQL
        \end{onecolentry}

        \vspace{0.2 cm}

        \begin{onecolentry}
            \textbf{Mensageria:} Apache Kafka, AWS SQS e Azure Service Bus
        \end{onecolentry}

        \vspace{0.2 cm}

        \begin{onecolentry}
            \textbf{Infrastrutura:} Terraform, Azure Biceps, AWS CDK, Docker, Kubernetes, AWS Lambda e Azure Functions
        \end{onecolentry}
        
\end{document}